% ============================================================
% 01-introduction.tex
% ============================================================
\section{Introduction}
\label{sec:intro}

Can Tho University (CTU) serves nearly 49,000 undergraduate students in the 2025--2026 academic year~\cite{ctu2025opening}. Its scale and evolving regulations create advisory needs across academic and administrative areas. CAAS supports chat- and voice-based admission counseling for information-technology programs~\cite{ma2023caas}, while REBot addresses academic-regulation advising through dense retrieval and category-guided graph reasoning~\cite{ma2025rebot}. Their reported evaluations support task-specific retrieval and processing, but each centers on an individual advisory domain, motivating a coordinated multi-domain architecture.

University counseling also spans heterogeneous knowledge forms: (1)~\textbf{academic curricula} encode hierarchies and prerequisites; (2)~\textbf{tuition schedules} vary by cohort and field; (3)~\textbf{exemption policies} and (4)~\textbf{scholarship policies} combine eligibility conditions with rates and thresholds; and (5)~\textbf{administrative regulations} comprise lengthy narrative documents. A single representation or retrieval method may therefore be unsuitable across all domains, motivating representation-specific evidence acquisition.

RAG and GraphRAG augment LLMs with documents and knowledge graphs~\cite{lewis2020rag,gao2023rag_survey,edge2024graphrag}, while agentic architectures support tool use and multi-agent coordination~\cite{yao2023react,wu2023autogen,schick2023toolformer}. These approaches do not prescribe how heterogeneous institutional evidence should be represented, routed, and processed. Known limitations in language-model arithmetic~\cite{qian2023limitations} further motivate deterministic calculation for rule-based financial queries.

Together with the representative literature synthesized in Section~2, these challenges motivate three gaps examined here: (1) limited evidence on bounded specialization under tool ambiguity; (2) limitations of uniform representations for heterogeneous knowledge; and (3) limited evidence on routing across isolated, representation-specific evidence paths.

The objective of this study is to develop and evaluate CTU-Chat for heterogeneous university counseling. To address the identified gaps, we formulate three research questions:
\begin{itemize}
    \setlength{\itemsep}{1.5pt}
    \setlength{\parsep}{0pt}
    \setlength{\parskip}{0pt}

    \item \textbf{RQ1 (Domain-Specialized Tool Disambiguation under Registry Expansion):}
    How do full-registry and Top-$k$ single-agent baselines compare with supervisor-routed generic and specialist gates in tool-selection reliability, decision-level end-to-end pass rate, and computational cost as semantically overlapping within-domain tools expand the registry from 11 to 51 tools?

    \item \textbf{RQ2 (Heterogeneous Knowledge Allocation \& Modular Ablations):}
    What differences in context and answer quality are observed between the full heterogeneous configuration, uniform retrieval baselines, and its component ablations?

    \item \textbf{RQ3 (Representation-Aware Retrieval Effectiveness):}
    How does the complete representation-aware evidence configuration compare with sparse, dense, hybrid, and reranked baselines in retrieval effectiveness, both overall and across query-complexity tiers?
\end{itemize}

We hypothesize that domain-scoped candidate retrieval and contrastive specialist policies can reduce confusion among semantically adjacent tools relative to a global Top-$k$ single-agent gate. This hypothesis is limited to the evaluated registry expansions and does not predict multi-agent dominance over a full-registry monolithic baseline on every accuracy or efficiency measure. CTU-Chat coordinates four specialists for academic, financial, scholarship, and general administrative inquiries, combining graph retrieval, structured access, deterministic calculation, and hybrid document retrieval. We evaluate it through retrieval stacking, end-to-end generation and ablation, and controlled tool-ambiguity stress tests against matched single-agent baselines.

Specifically, this work provides three primary contributions:
\begin{itemize}
    \setlength{\itemsep}{1.5pt}
    \setlength{\parsep}{0pt}
    \setlength{\parskip}{0pt}
    \item \textbf{C1 (Domain-Specialized Multi-Agent Architecture):} We implement a centralized supervisor coordinating four bounded specialists with isolated evidence representations, deterministic route repair, and partitioned tool privileges, and evaluate registries of up to 51 tools.
    \item \textbf{C2 (Heterogeneous Knowledge Allocation):} We model relational curricula as property graphs, bind statutory fee and exemption lookup to deterministic arithmetic, and index narrative regulations through parent-resolved document retrieval.
    \item \textbf{C3 (Tool-Ambiguity Stress Test and End-to-End Evaluation):} We evaluate semantic-neighbor ambiguity across 1,960 decisions, retrieval stacking across five configurations, and generation quality across seven configurations ($2,100$ turns) on a redesigned 100-query held-out benchmark.
\end{itemize}

The remainder of this paper is organized as follows. \textbf{Section~2} reviews related work and summarizes the research gaps. \textbf{Section~3} presents the CTU-Chat architecture, evidence-acquisition methods, and domain-specialized agents. \textbf{Section~4} reports the experimental setup, results, and discussion. Finally, \textbf{Section~5} concludes the paper and outlines future work.
